% !TEX encoding = UTF-8 Unicode
\documentclass[a4paper]{article}

\usepackage[T1]{fontenc}     % För svenska bokstäver
\usepackage[utf8]{inputenc}  % Teckenkodning UTF8
\usepackage[swedish]{babel}  % För svensk avstavning och svenska
                             % rubriker (t ex "Innehållsförteckning")
\usepackage{fancyvrb}        % För programlistor med tabulatorer
\fvset{tabsize=4}            % Tabulatorpositioner
\fvset{fontsize=\small}      % Lagom storlek för programlistor

\usepackage{graphicx}

\title{Dokumentnamn}
\author{Nils Nilsson}
\date{1 augusti 1994}        % Blir dagens datum om det utelämnas

\begin{document}             % Början på dokumentet

\maketitle                   % Skriver ut rubriken som vi
                             % definierade ovan med \title, \author
                             % och eventuellt \date


MATTIAS \--\-- ÄR  -{}- BÄST PÅ -\-- ALLT! \--

Här skriver jag något \emph{viktigt}. Och i
Java har vi använt klassen \texttt{Square}.

\section{Specialtecken}

Ej matematisk formel, utan \$7.



DoD-kursen pågår under vecka
1--3 av läsperiod ht1. Tyvärr 
är den inte längre \ldots
Telefon: 046--222 80 38. Dagens
datum: 6 september 2021.

Random mening --- fördtydligande --- mer text 


Min text\footnote{det är inte min egentligen?}\footnote{någon annans?} är skirevn här!


\section{Listor}


\begin{itemize}
	\item Min första punkt!
	\item Yay, detta är ju enkelt!
\end{itemize}

My top 3 word
\begin{enumerate}
	\item I
	\item AM
	\item BEST
\end{enumerate}

\begin{description}
	\item [SimpleWindow] Ett enkelt fönster
	\item [Scanner] Läsa text
	\item [Random] Slump!
\end{description}


Lite text och \\ lite mer
I denna text refererar jag till tabell~\ref{tab4}

\bigskip

Lite mer text, men som ligger en bit längre ned!

\begin{table}
\begin{tabular}{lllrr}
Rubrik1 & Rubrik2 & Rubrik3 & Rubrik4&      Rubrik5 \\
\hline
value&value & value & value & value \\
value & value & value & value & value  
\end{tabular}
\caption{Min tabelltext}
\label{tab4}
\end{table}

\section{Matematik}

Formeln $x = 3y - 2$ står inne i texten.

Däremot står
\begin{displaymath}
	x=3y-2
\end{displaymath}

\begin{equation}
	x=3y-2
	\label{eq1}
\end{equation}
\begin{equation}
	x=3y-2
	\label{eq2}
\end{equation}
\begin{equation}
	X_{pk} = X_{pk-1} + 4
	\label{eq3}
\end{equation}


\begin{equation}
	e^x = 1 + x + \frac{x^2}{2!} + \frac{x^3}{3!} + \cdots
\end{equation}

\begin{equation}
	\sin
\end{equation}

Titta i ekvation

\begin{figure}
\includegraphics[width=2cm]{../images/bild}
\caption{Min hejhopp-bild!}
\label{hejhopp}
\end{figure}

Referens till figur~\ref{hejhopp}

\end{document}               % Slut på dokumentet
